\section*{第 \ref{ch:b2:solid-liquid-diagrams} 章——二元固液相图}

\begin{solution}{exo:b2:solid-liquid-diagrams:1}
在 \qty{1300}{\degreeCelsius} 时，液相线位于 0.541，固相线位于 0.609。
0.40：仅有液相。0.58：液相（0.541）与固溶体（0.609）共存。0.70：
仅有固溶体。
\end{solution}

\begin{solution}{exo:b2:solid-liquid-diagrams:2}
固相分数 $(0.58 - 0.541)/(0.609 - 0.541) = 0.57$：固相 \qty{1.15}{mol}，
液相 \qty{0.85}{mol}。镍：固相中 $1.15 \times 0.609 = \qty{0.70}{mol}$，
液相中 $0.85 \times 0.541 = \qty{0.46}{mol}$，总量 \qty{1.16}{mol}，
$= 2.00 \times 0.58$。
\end{solution}

\begin{solution}{exo:b2:solid-liquid-diagrams:3}
恒压下 $v = n + 1 - \varphi$（无化学反应）。(a) $1 + 1 - 2 = 0$。
(b) $2 + 1 - 2 = 1$。(c) 液相与两种固溶体：$2 + 1 - 3 = 0$。
\end{solution}

\begin{solution}{exo:b2:solid-liquid-diagrams:4}
温度平稳降至 \qty{232}{\degreeCelsius}，锡凝固时出现水平平台，
随后再次平稳下降，在室温附近下降得更慢。平台上液相与固相在恒压下共存，
自由度为 0：
放出的热量由结晶过程提供，
在最后一部分液体凝固之前，温度不能改变。
\end{solution}

\begin{solution}{exo:b2:solid-liquid-diagrams:5}
苯：$\ln x_B = -(9870/8.314)(1/269.6 - 1/278.64) = -0.1429$，$x_B =
0.867$。萘：$\ln x_N = -(19\,100/8.314)(1/269.6 - 1/353.2) = -2.017$，
$x_N = 0.133$。两者之和为 1.000：\qty{269.6}{K}（$\qty{-3.5}{\degreeCelsius}$）为
低共熔温度，低共熔液相中萘的摩尔分数为 0.133。
\end{solution}

\begin{solution}{exo:b2:solid-liquid-diagrams:6}
液相冷却至铅一侧的液相线，此时开始析出 (Pb) 晶体；
液相中锡含量不断增加，直至 \qty{182.2}{\degreeCelsius} 处的低共熔点，
对应组成 62.1~\%。剩余液相在该处凝固为低共熔混合物。
低共熔混合物的质量分数为 $(30 - 18.91)/(62.13 - 18.91) = 0.26$。继续冷却时，
锡在 (Pb) 中的溶解度下降（虚线），少量 (Sn)
于 (Pb) 晶体内部析出。
\end{solution}

\begin{solution}{exo:b2:solid-liquid-diagrams:7}
按质量计 NaCl 为 23.3~\%：$23.3/76.7 = \qty{0.304}{kg}$，
即每千克水中含盐。在 \qty{-25}{\degreeCelsius} 下，低于低共熔温度，
无论比例如何都不可能存在液相：盐与冰以固态并存。
\end{solution}

\begin{solution}{exo:b2:solid-liquid-diagrams:8}
$x_{\mathrm B} = 0.5$ 位于 E$_1$ 与化合物之间：$\mathrm{AB_2}$
先结晶，液相向 E$_1$ 移动，在该处凝固为
A 与 $\mathrm{AB_2}$ 的低共熔混合物；最终为 A + $\mathrm{AB_2}$。
$x_{\mathrm B} = 0.9$ 位于 E$_2$ 与 B 之间：B 先结晶，液相
到达 E$_2$；最终为 $\mathrm{AB_2}$ + B。
\end{solution}

\begin{solution}{exo:b2:solid-liquid-diagrams:9}
在熔区后缘，液体凝固成的固体所含铜的摩尔分数
仅为液相中的 0.78 倍：被排出的铜留在液相中，
液相向前移动并逐渐富集，因此铜被推向棒的末端。
比值接近 1 时，每次通过只能去除杂质的一小部分，
需要多次通过，每次都从上一次留下的棒材开始。
\end{solution}

\begin{solution}{exo:b2:solid-liquid-diagrams:10}
推导同正文。对于稀液体，$\ln x_A \approx -x_B$，$1/T -
1/T^* \approx -\Delta T/T^{*2}$，因此 $x_B = \Delta_{\text{fus}}H_A\Delta T/(RT^{*2})$；
由 $x_B \approx n_B/n_A = bM_A$，得 $\Delta T = (RT^{*2}M_A/\Delta_{\text{fus}}H_A)\,b$。
苯：$K_f = 8.314 \times 278.64^2 \times 0.07811/9870 = \qty{5.11}{K.kg/mol}$。
\end{solution}

\begin{solution}{exo:b2:solid-liquid-diagrams:11}
每 \qty{100}{g}：镁 $19.0/24.31 = \qty{0.782}{mol}$，铅 $81.0/207.2 =
\qty{0.391}{mol}$；比值 $2.00$：$\ce{Mg2Pb}$。
\end{solution}

\begin{solution}{exo:b2:solid-liquid-diagrams:12}
平台长度与低共熔液相的质量成正比：未知合金中该液相占
$120/300 = 0.40$。在铅一侧，$w = 18.91 + 0.40 \times (62.13 -
18.91) = 36.2~\%$ 锡；在锡一侧，$w = 96.59 - 0.40 \times (96.59 - 62.13) =
82.8~\%$。冷却曲线的第一个转折点（对 36~\% 合金更高，
其液相线位于陡峭的铅一侧）或显微照片（初生 (Pb) 或初生 (Sn) 晶体）
可以区分二者。
\end{solution}

\begin{solution}{pb:b2:solid-liquid-diagrams:1}
\textbf{1.} 铅 \qty{327.5}{\degreeCelsius}，锡 \qty{231.9}{\degreeCelsius}。
\textbf{2.} 每 \qty{100}{g}：锡 $62.13/118.71 = \qty{0.5234}{mol}$，
铅 $37.87/207.2 = \qty{0.1828}{mol}$；$x_{\mathrm{Sn}} = 0.5234/0.7062 =
0.741$。
\textbf{3.} 点 (0, 327.5)、(100, 231.9)、(62.13, 182.2)、(18.91, 182.2)、
(96.59, 182.2)。
\textbf{4.} 如图所示：液相；L + (Pb)；L + (Sn)；(Pb)；(Sn)；(Pb) + (Sn)。
\textbf{5.} L（62.1~\% Sn）$\longrightarrow$ (Pb)（18.9~\%）+ (Sn)（96.6~\%）。
\textbf{6.} $v = 2 + 1 - 3 = 0$：温度与三个组成
均固定。
\textbf{7.} 富铅固溶体 (Pb)：铅溶解锡，因此
与液相平衡的固相中已含锡，比例可由
固相线读出。
\textbf{8.} 沿液相线向低共熔点移动，液相中锡含量增加，
直至 62.13~\%。
\textbf{9.} 62.13~\% 的液相与含锡 18.91~\% 的 (Pb)。
\textbf{10.} (Pb)：$(62.13 - 40)/(62.13 - 18.91) = 0.512$；液相 0.488。
\textbf{11.} 含锡 18.91~\% 的 (Pb) 与 96.59~\% 的 (Sn)；(Sn)：$(40 - 18.91)/(96.59 -
18.91) = 0.271$；(Pb)：0.729。
\textbf{12.} 低共熔组织占 48.8~\%（刚好在其上方存在的液相）；
初生 (Pb) 占 51.2~\%。第 11 题中的 (Pb) 是
初生晶体与低共熔组织中 (Pb) 片层之和。
\textbf{13.} 液相线处出现转折，随后下降变慢，然后在
\qty{182.2}{\degreeCelsius} 出现平台，持续时间为低共熔合金的 0.488 倍。
\textbf{14.} $\ln x_{\mathrm{Sn}} = -(7030/8.314)(1/T - 1/505.1)$。
\textbf{15.} $\ln x_{\mathrm{Bi}} = -(11\,300/8.314)(1/T - 1/544.6)$。
\textbf{16.} $x_{\mathrm{Sn}}(T) + x_{\mathrm{Bi}}(T) = 1$。
\textbf{17.} 在 \qty{110}{\degreeCelsius}：$0.587 + 0.349 = 0.936$；在
\qty{120}{\degreeCelsius}：$0.621 + 0.382 = 1.003$；在
\qty{130}{\degreeCelsius}：$0.655 + 0.417 = 1.072$。两者之和在
\qty{119.5}{\degreeCelsius} 附近达到 1：取整到度即 \qty{120}{\degreeCelsius}。
\textbf{18.} $x_{\mathrm{Bi}} = 0.38$；按质量计 $0.38 \times 208.98/(0.38 \times
208.98 + 0.62 \times 118.71) = 0.52$。
\textbf{19.} 模型给出 \qty{120}{\degreeCelsius} 与 52~\% 铋，而评估相图为
\qty{138.8}{\degreeCelsius} 与 57~\%；模型约
低 \qty{19}{\degreeCelsius}。
\textbf{20.} 析出的锡并非纯锡，而是含铋的溶液，
其化学势低于纯锡，因此固相更稳定，在相同液相组成下
它在更高温度下从液相中结晶。锡的液相线分支被抬高，
并在更高处与铋的分支相交
（非理想液体还会附加其自身的修正）。
\textbf{21.} 合金在单一温度下熔化和凝固，即体系的最低温度：
不存在糊状区间，因此移动的接头在此期间不会开裂。
\textbf{22.} 铅有毒，电子废弃物最终进入环境。
\textbf{23.} 对热敏元件可在较低温度下焊接；
但接头在使用中软化温度较低，且铋使其
变脆。
\textbf{24.} 铋 \qty{570}{g}，锡 \qty{430}{g}。
\textbf{25.} 理想液体模型预测低共熔点为 $\boldsymbol{T_E
\approx \qty{120}{\degreeCelsius}}$，
而评估值为 \qty{138.8}{\degreeCelsius}。
\end{solution}
